\documentclass[11pt]{article}
\usepackage[a4paper,margin=1in]{geometry}
\usepackage{amsmath,amssymb,amsthm}
\usepackage{array}
\usepackage{hyperref}

\newcommand{\dd}{\mathrm{d}}
\newcommand{\sgn}{\operatorname{sgn}}
\newcommand{\cA}{\mathcal{A}}
\newcommand{\cD}{\mathcal{D}}
\newcommand{\avg}[1]{\langle #1\rangle}
\newcommand{\vt}{v_t}
\renewcommand{\Re}{\operatorname{Re}}
\renewcommand{\Im}{\operatorname{Im}}

\newtheorem{theorem}{Theorem}
\newtheorem{lemma}[theorem]{Lemma}
\newtheorem{proposition}[theorem]{Proposition}
\theoremstyle{remark}
\newtheorem{remark}[theorem]{Remark}

\title{The linear heat-flux closure in three velocity dimensions\\
for a general conserving collision operator, with numbers for Landau}
\author{Working note}
\date{16 November 2026}

\begin{document}
\sloppy
\maketitle

\noindent\emph{Summary.} The previous notes left the Landau operator open
for two reasons: no closed form, and, in three dimensions, a coupling to the
viscous stress that breaks the positivity argument. This note gets as far as
the argument goes.
\begin{itemize}
  \item \textbf{Proved} (Theorem~\ref{thm:3d}): for any conserving,
    symmetric, non-negative linearized collision operator in three velocity
    dimensions, the quasi-static closure obeys
    $|\chi(k)/\kappa_{\mathrm F}-1|\le K\,k^2$ at small $k$, two-sided, with
    $K$ given explicitly by three Chapman--Enskog-type quantities. The
    Landau operator is covered, conditionally on the existence statement in
    Assumption~(S).
  \item \textbf{Proved} (Theorem~\ref{prop:K}): the exact coefficient of
    $k^2$ in the departure from Fourier's law, for any such operator, under
    a solvability assumption on two Chapman--Enskog problems. It gives
    $\tfrac{2}{15}$ for the Dougherty operator, and in one velocity
    dimension it shows that the constants of the earlier notes are exact.
  \item \textbf{Computed} for the like-particle Landau operator: the
    conductivity, the viscosity, the exact $k^2$ coefficient, and the
    closure itself up to $k\lambda\approx1$, with a polynomial Galerkin
    method validated against Braginskii's coefficients and the Dougherty
    closed form.
\end{itemize}
\textbf{Not proved for Landau:} positivity at all $k$, the inequality
$\chi<\kappa_{\mathrm F}$, and the collisionless limit. Assumptions (R) and (S) are proved in
\texttt{landau\_solvability.tex} and
\texttt{landau\_quasistatic\_limit.tex}, from published coercivity
estimates. The Landau numbers are numerical, not certified. Electron--ion
collisions are not included. No second reader has checked the proofs.

\section{Setting}

Units: density, mass and thermal speed $\sqrt{T_0/m}$ equal to one. $M$ is
the unit Maxwellian on $\mathbb{R}^3$, $\avg{\cdot}$ its average. All
functions are axisymmetric about the $x$ axis, and
$H=\operatorname{span}\{1,v_x,|v|^2\}$. Write
\[
  Q=\tfrac12v_x(|v|^2-5),\qquad
  \vartheta=v_\perp^2-2v_x^2,\qquad
  T=\tfrac13\avg{(|v|^2-3)h},\qquad q=\avg{Qh}.
\]

\paragraph{Assumption (A3).} $\cD(f,h)=\avg{f\Lambda h}$ is a real,
symmetric, non-negative bilinear form whose domain contains the polynomials,
invariant under $v_x\to-v_x$, with $\cD[f]:=\cD(f,f)=0$ exactly for $f\in H$.

For real $\sigma$ put
$\avg{\sigma,\Lambda^+\sigma}=\sup_f[2\avg{\sigma f}-\cD[f]]$. Then
$\avg{\sigma f}^2\le\avg{\sigma,\Lambda^+\sigma}\,\cD[f]$, and if
$\cD(f_0,\cdot)=\avg{\sigma\,\cdot}$ then
$\avg{\sigma,\Lambda^+\sigma}=\cD[f_0]$.

\paragraph{Assumption (S).} For the given $k\ne0$ there is $g=g_e+iw$, with
$g_e$ real and even in $v_x$, $w$ real and odd, in the form domain and with
finite polynomial moments, solving $ikv_x\,g+\Lambda g=1$ weakly; and the
quasi-static limit of the potential-driven response of
$\partial_th+v_x\partial_xh+\Lambda h=a\,v_x$ is given by its moments, as in
the earlier notes: $\hat q/\hat T\to q_1/T_1$ with $q_1=\avg{Qg}$,
$T_1=\tfrac13\avg{(|v|^2-3)g}$.

\medskip
(S) was proved for the BGK and Dougherty operators. For the Landau operator
with Coulomb interaction the linearized operator has no spectral gap.
Existence of a stationary solution for which all the conclusions below hold
is proved in \texttt{landau\_solvability.tex}, and the identification of
the quasi-static limit with it in \texttt{landau\_quasistatic\_limit.tex}.

\paragraph{Chapman--Enskog quantities.} Let $\Xi$ solve $\Lambda\Xi=Q$ with
$\avg{v_x\Xi}=0$. Define
\begin{equation}
  \kappa_{\mathrm F}=\avg{Q\,\Xi},\qquad
  \eta=\tfrac1{12}\avg{\vartheta,\Lambda^+\vartheta},\qquad
  \chi_*=1-v_x^2+\frac{v_x\,\Xi}{\kappa_{\mathrm F}},\qquad
  C_*=\avg{\chi_*,\Lambda^+\chi_*}.
\end{equation}
$\kappa_{\mathrm F}$ is the heat conductivity and $\eta$ the viscosity
($\pi_{xx}=-\tfrac43\eta\,\partial_xu$). For the Dougherty operator
$\kappa_{\mathrm F}=5/(6\nu)$ and $\eta=1/(2\nu)$.

\section{Theorem}

\begin{theorem}\label{thm:3d}
Assume (A3) and (S), and let $n_1=\avg{g}$,
$X=\avg{\vartheta\,g_e}$, $X_{\max}=(12\,\eta\,C_*)^{1/2}$.
\begin{enumerate}
  \item[(i)] $q_1=i/k$, $\avg{v_x^2g}=0$, and
    $T_1=-n_1+\tfrac13X$.
  \item[(ii)] $\displaystyle\frac{1}{k^2\kappa_{\mathrm F}}\le n_1\le
    \frac{1}{k^2\kappa_{\mathrm F}}+C_*$ and $|X|\le X_{\max}$.
  \item[(iii)] If $\kappa_{\mathrm F}X_{\max}k^2<3$, then $T_1<0$, the
    closure is $\hat q=-i\,\sgn(k)\,\theta\,\hat T$ with
    $\theta=-1/(|k|T_1)>0$, and $\chi=\theta/|k|$ satisfies
    \begin{equation}
      \frac{\kappa_{\mathrm F}}
           {1+\kappa_{\mathrm F}\bigl(C_*+\tfrac13X_{\max}\bigr)k^2}
      \;\le\;\chi(k)\;\le\;
      \frac{\kappa_{\mathrm F}}{1-\tfrac13\kappa_{\mathrm F}X_{\max}k^2}.
      \label{eq:bounds}
    \end{equation}
\end{enumerate}
\end{theorem}

\begin{proof}
Take $k>0$.

\emph{(i).} Testing the equation with $1$, $v_x$, $|v|^2$, which lie in the
kernel of the form: $ik\avg{v_xg}=1$, $ik\avg{v_x^2g}=0$,
$ik\avg{v_x|v|^2g}=\avg{|v|^2}=3$. Hence $q_1=(3-5)/(2ik)=i/k$. Also
$\avg{(|v|^2-3)g}=\avg{v_\perp^2g}-3n_1$ and
$\avg{\vartheta g}=\avg{v_\perp^2g}$, which gives $3T_1=-3n_1+X$; and $X$ is
real with $\avg{\vartheta g}=\avg{\vartheta g_e}$ because $\vartheta$ is
even.

\emph{Energy identity.} The real and imaginary parts of the equation are
\begin{equation}
  -kv_xw+\Lambda g_e=1,\qquad kv_xg_e+\Lambda w=0 .
  \label{eq:sys}
\end{equation}
Testing the first with $g_e$ and the second with $w$ and subtracting,
\begin{equation}
  n_1=\cD[g_e]+\cD[w].
  \label{eq:energy}
\end{equation}

\emph{Lower bound on $\cD[w]$.} Since $Q$ is odd, $q_1=i\avg{Qw}$, so
$\avg{Qw}=1/k$. Then
$k^{-2}=\avg{Qw}^2\le\avg{Q,\Lambda^+Q}\,\cD[w]=\kappa_{\mathrm F}\cD[w]$.
With \eqref{eq:energy} this is the lower bound in (ii).

\emph{Upper bound on $n_1$.} Let
$\psi_*=(v_x-\Xi/\kappa_{\mathrm F})/k$. Then
$1-kv_x\psi_*=\chi_*$, $\cD[\psi_*]=\cD[\Xi]/(k\kappa_{\mathrm F})^2
=1/(k^2\kappa_{\mathrm F})$, and $\chi_*$ is orthogonal to $1$ and $|v|^2$:
$\avg{\chi_*}=1-1+0$ and
$\avg{|v|^2\chi_*}=3-5+\avg{(2Q+5v_x)\Xi}/\kappa_{\mathrm F}=0$. Write
\[
  n_1=\avg{g_e}=\avg{g_e\chi_*}+k\avg{v_xg_e\,\psi_*}.
\]
The first term is at most $C_*^{1/2}\cD[g_e]^{1/2}$. For the second, test
the second equation of \eqref{eq:sys} with $\psi_*$:
$k\avg{v_xg_e\psi_*}=-\cD(w,\psi_*)\le\cD[w]^{1/2}(k^2\kappa_{\mathrm F})^{-1/2}$.
By the Cauchy--Schwarz inequality in $\mathbb{R}^2$ and \eqref{eq:energy},
$n_1\le(C_*+1/(k^2\kappa_{\mathrm F}))^{1/2}n_1^{1/2}$.

\emph{Bound on $X$.} From the two bounds, $\cD[g_e]=n_1-\cD[w]\le C_*$. Since
$\vartheta$ is orthogonal to $H$,
$X^2\le\avg{\vartheta,\Lambda^+\vartheta}\,\cD[g_e]\le12\,\eta\,C_*$.

\emph{(iii).} $-T_1=n_1-\tfrac13X$ lies between
$1/(k^2\kappa_{\mathrm F})-\tfrac13X_{\max}$ and
$1/(k^2\kappa_{\mathrm F})+C_*+\tfrac13X_{\max}$. Under the stated condition
the lower end is positive, $q_1/T_1=-i\theta$ with $\theta=-1/(kT_1)$, and
$\chi=1/(k^2(-T_1))$ gives \eqref{eq:bounds}.
\end{proof}

\begin{remark}
In one velocity dimension $\avg{v^2g}=0$ makes $T_1=-n_1$, a diagonal
resolvent element, and positivity holds at every $k$. In three dimensions
$T_1$ contains the cross term $X$, which is the perpendicular temperature
response; the argument controls it only by $X_{\max}$, so positivity and the
upper bound are obtained only for $k^2<3/(\kappa_{\mathrm F}X_{\max})$. For
the Dougherty operator both hold at every $k$, by the closed form.
\end{remark}

\section{The exact coefficient of $k^2$}

\paragraph{Assumption (R).} The following second-order Chapman--Enskog
problems have solutions in the form domain: $\Lambda G=\chi_*$, and
$\Lambda\Upsilon=v_x\tilde G$, where $\tilde G=G+h$ with
$h\in\operatorname{span}\{1,|v|^2\}$ chosen so that
$\avg{\tilde G}=\avg{v_x^2\tilde G}=0$. Put
\begin{equation}
  X_0=\avg{\vartheta\,G}=\avg{\vartheta,\Lambda^+\chi_*},\qquad
  B=\cD[\Upsilon]=\avg{v_x\tilde G,\Lambda^+v_x\tilde G}.
\end{equation}
For the BGK and Dougherty operators all of these are polynomials. For the
Landau operator (R) is proved in \texttt{landau\_solvability.tex}, from
published coercivity estimates.

\begin{theorem}[Leading coefficient]\label{prop:K}
Assume (A3), (S) and (R), and write $\varepsilon=k^2\kappa_{\mathrm F}$.
Then
\begin{align}
  \frac1\varepsilon+C_*-B\,k^2\;&\le\;n_1\;\le\;\frac1\varepsilon+C_*,
  \label{eq:n1two}\\
  |X-t\varepsilon\,X_0|\;&\le\;(12\,\eta\,B)^{1/2}\,|k|,\qquad
  t\varepsilon=\bigl(1-\varepsilon C_*+\varepsilon k^2B\bigr)^{-1},
  \label{eq:Xclose}
\end{align}
and consequently, as $k\to0$,
\begin{equation}
  1-\frac{\chi(k)}{\kappa_{\mathrm F}}=K_{\mathrm{ex}}\,k^2+O(|k|^3),\qquad
  K_{\mathrm{ex}}=\kappa_{\mathrm F}\bigl(C_*-\tfrac13X_0\bigr)
  =\kappa_{\mathrm F}\,
  \bigl\langle\chi_*-\tfrac13\vartheta,\;\Lambda^+\chi_*\bigr\rangle .
\end{equation}
\end{theorem}

\begin{proof}
For even $\varphi$ with $\avg{v_x^2\varphi}=0$ such that
$\Lambda f_\varphi=v_x\varphi$ is solvable in the form domain, define
\[
  \mathcal{Q}(\varphi,\varphi')=\cD(\varphi,\varphi')+k^2\cD(f_\varphi,f_{\varphi'}),
  \qquad \mathcal{Q}[\varphi]=\mathcal{Q}(\varphi,\varphi).
\]
By the second equation of \eqref{eq:sys}, $f_{g_e}=-w/k$.

\emph{Step 1: $\mathcal{Q}(g_e,\varphi)=\avg{\varphi}$ and
$\mathcal{Q}[g_e]=n_1$.} Testing the first equation of \eqref{eq:sys} with
$\varphi$ gives $\cD(g_e,\varphi)=\avg{\varphi}+k\avg{v_xw\,\varphi}$.
Testing $\Lambda f_\varphi=v_x\varphi$ with $w$ gives
$k^2\cD(-w/k,f_\varphi)=-k\avg{v_xw\,\varphi}$. Adding gives the first
claim; the second is \eqref{eq:energy}. Hence, for every such $\varphi$,
\begin{equation}
  \mathcal{Q}[g_e-\varphi]=n_1-\bigl(2\avg{\varphi}-\mathcal{Q}[\varphi]\bigr).
  \label{eq:gap}
\end{equation}

\emph{Step 2: a trial function.} Let $\varphi_0=\tfrac12(5-|v|^2)$, so
$\avg{\varphi_0}=1$, $\avg{v_x^2\varphi_0}=0$, $\cD[\varphi_0]=0$,
$v_x\varphi_0=-Q$ and $f_{\varphi_0}=-\Xi$. Let
$\varphi'=\varphi_0+\varepsilon\tilde G$; then $\avg{\varphi'}=1$,
$f_{\varphi'}=-\Xi+\varepsilon\Upsilon$, and
\[
  \mathcal{Q}[\varphi']=\varepsilon^2\cD[G]
  +k^2\bigl(\cD[\Xi]-2\varepsilon\,\cD(\Xi,\Upsilon)+\varepsilon^2B\bigr).
\]
Here $\cD[G]=C_*$, $\cD[\Xi]=\kappa_{\mathrm F}$, and
$\cD(\Xi,\Upsilon)=\avg{v_x\Xi\,\tilde G}$. Since
$v_x\Xi=\kappa_{\mathrm F}(\chi_*-1+v_x^2)$ and $\chi_*$ is orthogonal to
$h$,
$\avg{v_x\Xi\,\tilde G}=\kappa_{\mathrm F}\bigl(\avg{\chi_*G}-\avg{\tilde G}
+\avg{v_x^2\tilde G}\bigr)=\kappa_{\mathrm F}C_*$. With
$k^2\kappa_{\mathrm F}=\varepsilon$,
\begin{equation}
  \mathcal{Q}[\varphi']=\varepsilon\bigl(1-\varepsilon C_*+\varepsilon k^2B\bigr).
\end{equation}

\emph{Step 3: \eqref{eq:n1two}.} The upper bound is
Theorem~\ref{thm:3d}(ii). Take $\varphi=t\varphi'$ in \eqref{eq:gap} with
$t=1/\mathcal{Q}[\varphi']$, so that
$2\avg{\varphi}-\mathcal{Q}[\varphi]=1/\mathcal{Q}[\varphi']$. Since the
left side of \eqref{eq:gap} is non-negative and $1/(1-x)\ge1+x$,
\[
  n_1\ge\frac{1}{\mathcal{Q}[\varphi']}
  \ge\frac1\varepsilon\bigl(1+\varepsilon C_*-\varepsilon k^2B\bigr)
  =\frac1\varepsilon+C_*-Bk^2
\]
for $k$ small enough that $\mathcal{Q}[\varphi']>0$ and
$\varepsilon(C_*-k^2B)<1$.

\emph{Step 4: \eqref{eq:Xclose}.} By \eqref{eq:gap} and Step 3,
$\cD[g_e-t\varphi']\le\mathcal{Q}[g_e-t\varphi']
=n_1-1/\mathcal{Q}[\varphi']\le Bk^2$. Since $\vartheta$ is orthogonal to
$H$, $\avg{\vartheta\,t\varphi'}=t\varepsilon\avg{\vartheta\tilde G}
=t\varepsilon X_0$, and
$|\avg{\vartheta(g_e-t\varphi')}|\le(12\eta)^{1/2}\cD[g_e-t\varphi']^{1/2}$.

\emph{Step 5.} By Theorem~\ref{thm:3d}(i),
$-T_1=n_1-\tfrac13X=\varepsilon^{-1}+C_*-\tfrac13X_0+O(|k|)$, using
$t\varepsilon=1+O(k^2)$. Hence
$\chi/\kappa_{\mathrm F}=1/(\varepsilon(-T_1))
=1-\varepsilon(C_*-\tfrac13X_0)+O(|k|^3)$.
\end{proof}

\begin{remark}[One velocity dimension]
The same proof, with $\varphi_0=\tfrac12(3-v^2)$,
$\varepsilon=\tfrac12k^2\kappa_{\mathrm F}$ and no cross term, gives
$2/(k^2\kappa_{\mathrm F})+C_*-Bk^2\le n_1\le2/(k^2\kappa_{\mathrm F})+C_*$
and therefore
$1-\chi/\kappa_{\mathrm F}=\tfrac12C_*\kappa_{\mathrm F}\,k^2+O(k^4)$. So the
constant $K_2$ of the earlier notes is the exact coefficient, which was
there only observed for velocity-dependent $\nu$.
\end{remark}

\begin{remark}
The error term in three dimensions is stated as $O(|k|^3)$ because $X$ is
controlled only to $O(|k|)$. Numerically the error is $O(k^4)$, with
coefficient close to $-\kappa_{\mathrm F}B$ for the Landau operator.
\end{remark}

\section{Numerical method and validation}
\label{sec:num}

Scratch computations on 16 November 2026, not part of the repository.

\paragraph{Operator.} The like-particle Landau operator, linearized, in the
weak form
\[
  \cD(g,h)=\frac{3\sqrt\pi}{4\,\tau}\iint M(v)M(v')\,
  (\nabla g-\nabla'g')\cdot U(v-v')\,(\nabla h-\nabla'h')\,\dd v\,\dd v',
\]
\[
  U(w)=\frac{|w|^2I-ww^{\mathsf T}}{|w|^3},
\]
with $\tau$ Braginskii's like-particle collision time and $\lambda=\tau$ in
these units. This is one species colliding with itself.

\paragraph{Discretization.} Galerkin in the polynomials
$\psi_n(v_x)L_m(v_\perp^2/2)$ of total degree $n+2m\le d$. In the variables
$(v\pm v')/\sqrt2$ the integrand is a polynomial times $1/|u|$, and the
matrix is computed by Gauss quadrature that is exact for it.

\paragraph{Validation.}
\begin{itemize}
  \item $\Lambda$ annihilates $1$, $v_x$, $|v|^2$ to $10^{-14}$, is
    symmetric to $10^{-15}$ and has no negative eigenvalue beyond rounding.
  \item $\kappa_{\mathrm F}=3.90625\,\tau$ at $d=6$, which is Braginskii's
    $3.906$, and converges to $3.9502\,\tau$;
    $\eta=0.9653\,\tau$ against Braginskii's $0.96$.
  \item With the Dougherty matrix in place of the Landau one, the same code
    returns $\kappa_{\mathrm F}=5/6$, $\eta=1/2$ and
    $K_{\mathrm{ex}}=0.133333=\tfrac{2}{15}$, and the Galerkin closure agrees
    with the closed form of the previous note to six digits for
    $k\lambda\le2$.
  \item The estimates \eqref{eq:n1two} and \eqref{eq:Xclose} hold at
    $k\lambda=0.003$, $0.01$, $0.03$, $0.1$ for the Dougherty operator
    ($B=0.365$) and the Landau operator ($B\approx578\,\tau^3$), except that
    for Dougherty at $k\lambda=0.003$ the upper bound in \eqref{eq:n1two} is
    exceeded by $10^{-5}$, which is rounding in the subtraction of
    $1/\varepsilon\approx10^5$.
\end{itemize}

\section{Results for the Landau operator}

\begin{center}
\renewcommand{\arraystretch}{1.2}
\begin{tabular}{rrrrrr}
  $d$ & $\kappa_{\mathrm F}/\tau$ & $\eta/\tau$ & $C_*$
      & $\avg{\vartheta,\Lambda^+\chi_*}$ & $K_{\mathrm{ex}}/\tau^2$ \\ \hline
  6  & 3.9063 & 0.9653 & 4.308 & 0.883 & 15.68 \\
  8  & 3.9501 & 0.9653 & 5.958 & 0.813 & 22.46 \\
  10 & 3.9502 & 0.9653 & 5.937 & 0.815 & 22.38 \\
  12 & 3.9502 & 0.9653 & 5.928 & 0.815 & 22.34 \\
  16 & 3.9502 & 0.9653 & 5.929 & 0.815 & 22.35 \\
\end{tabular}
\end{center}
So $1-\chi/\kappa_{\mathrm F}\approx22.35\,(k\lambda)^2$, and the rigorous
constant in \eqref{eq:bounds} is
$\kappa_{\mathrm F}(C_*+\tfrac13X_{\max})=34.3\,\tau^2$ with
$X_{\max}=8.29$; the upper bound applies for $k\lambda<0.30$.

\begin{center}
\renewcommand{\arraystretch}{1.2}
\begin{tabular}{rrrrrr}
  $k\lambda$ & $\chi/\kappa_{\mathrm F}$ ($d=12$) & ($d=16$)
     & $1-22.35(k\lambda)^2$ & lower bound & upper bound \\ \hline
  0.005 & 0.99944 & 0.99944 & 0.99944 & 0.99914 & 1.00027 \\
  0.01  & 0.99779 & 0.99779 & 0.99777 & 0.99658 & 1.00109 \\
  0.02  & 0.99140 & 0.99140 & 0.99106 & 0.98645 & 1.00438 \\
  0.05  & 0.95426 & 0.95429 & 0.94413 & 0.92095 & 1.02805 \\
  0.1   & 0.87095 & 0.87112 & 0.77652 & 0.74442 & 1.12249 \\
  0.2   & 0.72711 & 0.72717 & ---     & 0.42135 & 1.77459 \\
  0.3   & 0.62603 & 0.62758 & ---     & 0.24450 & --- \\
  0.5   & 0.50159 & 0.50300 & ---     & 0.10435 & --- \\
  1.0   & 0.34975 & 0.35057 & ---     & 0.02830 & --- \\
\end{tabular}
\end{center}
\begin{itemize}
  \item The bounds of Theorem~\ref{thm:3d} hold at every point, as do the
    two inequalities of part (ii).
  \item The closure is converged in $d$ to four digits for
    $k\lambda\le0.2$ and to two or three for $k\lambda$ up to $1$. Beyond
    that a polynomial basis is not adequate, and the collisionless regime is
    not reached.
  \item $\chi<\kappa_{\mathrm F}$ at every computed point.
\end{itemize}

\section{Comparison across operators}

To compare operators with different collision-time conventions, normalize by
the conductivity: with $\vt=1$, write
$1-\chi/\kappa_{\mathrm F}\approx c\,(k\,\kappa_{\mathrm F})^2$.
\begin{center}
\renewcommand{\arraystretch}{1.25}
\begin{tabular}{llp{5.2cm}}
  Operator & $c = K/\kappa_{\mathrm F}^2$ & Status of $K$ \\ \hline
  BGK, constant $\nu$, 1D & $0.444$ & proved, sharp \\
  Dougherty, 1D & $0.333$ & proved, sharp \\
  BGK, $\nu_0/(1+|v|^3)$, 1D & $6.78$ & proved, sharp \\
  Dougherty, 3D & $0.192$ & proved (asymptotic) \\
  Landau, like-particle, 3D & $1.43$ & proved formula; numerical value \\
\end{tabular}
\end{center}
At equal conductivity the Landau operator departs from Fourier's law about
$7.5$ times faster in $k^2$ than the Dougherty operator: a 13\% deficit at
$k\lambda=0.1$, where $\lambda=\vt\tau$, against about 3\% for Dougherty at
the same $k\kappa_{\mathrm F}$. This is the like-particle operator. With
electron--ion scattering added the literature reports a much larger
coefficient; that case is not computed here.

A harmonic-mean interpolant built from $\kappa_{\mathrm F}$ and the
conjectured collisionless value $3\sqrt{2/\pi}$ is within 2\% of the
computed closure at $k\lambda=0.1$ and above it by 3\%, 9\% and 8\% at
$k\lambda=0.2$, $0.5$ and $1$. The sign is the one reported for the Landau
operator, and opposite to the Dougherty case. Given the convergence at
$k\lambda\ge0.5$ and the unproved limit, treat this as indicative.

\section{What is not done}

\begin{enumerate}
  \item \textbf{Assumption (S) for Landau.} Existence of the stationary
    solution is proved in the note \texttt{landau\_solvability.tex}, and
    the quasi-static limit in \texttt{landau\_quasistatic\_limit.tex}.
    Existence of the time-dependent weak solution in the weighted class used
    there is taken from a standard method and is not written out.
  \item \textbf{Positivity: improved range.} The note
    \texttt{landau\_positivity.tex} proves $\chi>0$ for
    $k^2\kappa_{\mathrm F}\eta<3$ ($k\lambda<0.89$ for Landau), reduces the
    question to the perpendicular temperature response, and records that
    positivity at all wavenumbers is still not proved for Landau.
  \item \textbf{All wavenumbers.} Positivity of $\chi$, and
    $\chi<\kappa_{\mathrm F}$, are proved only in the Dougherty case in
    three dimensions. For a general operator the theorem covers
    $k^2<3/(\kappa_{\mathrm F}X_{\max})$. A variational principle for the
    coupled heat-flux and stress response would be needed to go further.
  \item \textbf{Collisionless limit.} Not proved for Landau and not reached
    numerically.
  \item \textbf{Certification.} The Landau constants come from a Galerkin
    truncation. They are converged to the digits shown, but there is no
    error bound. Theorem~\ref{prop:K} is proved under (R), which holds
    for Landau by \texttt{landau\_solvability.tex}.
  \item \textbf{Physics.} One species, like-particle collisions, no
    electron--ion scattering, no magnetic field, linear, quasi-static,
    single drive.
  \item \textbf{Prior art.} Not searched. Wavenumber-dependent transport
    coefficients for the Landau operator have a large literature, and
    $K_{\mathrm{ex}}$ is a Burnett-type coefficient that may be tabulated.
\end{enumerate}

\section{Use in the benchmark}

\begin{itemize}
  \item \textbf{A test that separates ground truths.} The single number
    $c=K/\kappa_{\mathrm F}^2$ distinguishes the operators above by a factor
    of 35 between extremes. A closure trained on one and evaluated on
    another can be scored against it without a kinetic solve.
  \item \textbf{A certificate at small $k$.} For any conserving operator,
    \eqref{eq:bounds} brackets the exact closure using only Chapman--Enskog
    solves.
  \item \textbf{For the tournament.} The Dougherty ground truth has the
    smallest coefficient in the table. If the benchmark is meant to exercise
    departures from the local closure at moderate $k\lambda$, a Landau
    ground truth departs several times faster.
\end{itemize}

\end{document}
